% !TEX program = lualatex
% =============================================================
% OnePage LG — Hanaro / Antes / Depois / Resultado
%
%   BASEADO ESTRITAMENTE NO MODELO CORPORATIVO FORNECIDO.
%   Este arquivo altera somente o conteudo. Layout, cores, fontes,
%   coordenadas, icones e motor grafico permanecem em onepagelg.sty.
%
%   Compilar: lualatex onepage-lg.tex   (2x)
% =============================================================
\documentclass{article}
\usepackage{onepagelg}
\begin{document}

% ======================= CABECALHO E RODAPE =======================
\avisotopo{LGE — Uso Interno}
\titulo{1. Hanaro — Scrap \& IF Cost}
\subtitulo{Plataforma para centralizar dados de Material Scrap / IF Cost, ampliar a rastreabilidade e automatizar o fluxo de gestão de perdas.}
\colunas{Antes}{Depois}{Resultado}
\rodape{| \ C\,O\,N\,F\,I\,D\,E\,N\,C\,I\,A\,L\ |}{0 / 7}

% ======================= COLUNA 1 — ANTES =======================
\tituloantes{Contexto / Problema}
\textoantes{Dados de Scrap ainda dependem de consolidação manual e de fontes
  distribuídas, gerando retrabalho e divergências nos indicadores de produção.}

% Textos mantidos curtos para respeitar a densidade do modelo original.
\ponto[1.18]{ico-repetitivo.png}{A consolidação manual do Scrap exige esforço operacional recorrente.}
\ponto[1.08]{ico-sobrecarga.png}{Produção e Qualidade precisam cruzar informações antes da análise.}
\ponto[0.90]{ico-planilhas.png}{Regras de escopo e ajustes ainda dependem de planilhas.}

\alerta{ico-seguranca.png}{Baixa rastreabilidade e divergências entre fontes}

% quadro de baixo: titulo, o fluxo "de ... para ..." e os numeros de impacto
\tituloquadro{Situação atual}
\fluxo{De 1,5h a 2h/dia}{até 504 h/ano}
\destaque[0.96]{ico-pessoas.png}{Fluxo manual e fragmentado}
\destaque[0.96]{ico-oportunidades.png}{Sem base única para Scrap / IF Cost}

% ======================= COLUNA 2 — DEPOIS =======================
\titulodepois{Melhoria Esperada}
\textodepois{O Hanaro centraliza os dados de Material Scrap / IF Cost, padroniza a
  ingestão e conecta monitoramento, revisões, relatórios e planos de ação em um
  fluxo único e rastreável.}

\melhoria{Fluxo único, da ingestão à análise, revisão e ação.}
\melhoria{Base preparada para automação via RPA e PoCs controladas de IA.}

\abas{HANARO}{DADOS}

% \area{angulo}{rotulo} — o rotulo se posiciona sozinho pelo angulo
\area{90}{Ingestão}
\area{45}{Scrap}
\area{0}{IF Cost}
\area{-45}{Revisões}
\area{270}{Metas}
\area{225}{Relatórios}
\area{180}{Ações}
\area{135}{Alertas}

\afirmacao{Dashboard e Modo TV apoiam o monitoramento e a priorização.}
\afirmacao{Revisões, evidências e planos de ação ampliam a rastreabilidade.}
\afirmacao{RPA e IA local seguem como frentes de prospecção técnica / PoC, sem bloquear o MVP.}

\ilustracoes{robo-analista.png}{troca.png}{lg-ia.png}

% ======================= COLUNA 3 — RESULTADO =======================
\tituloresultado{Efeito da Melhoria}

\efeito{Criar uma visão única e confiável de Scrap / IF Cost}
\efeito{Ampliar a rastreabilidade da ocorrência até a ação}
\efeito{Reduzir a consolidação manual e acelerar decisões}

% \graficobarra{categoria}{rotulo do valor}{altura em cm}
\graficoeixo{Fluxo}
\graficobarra{Antes}{Manual}{1.04}
\graficobarra{Hanaro}{Integrado}{2.42}
\graficonota{MVP + governança}

\impacto{ico-areas.png}{Impacto no Negócio}{Produção / Qualidade}
\resultado{Meta M3}{Módulo funcional validado}

% =============================================================
\onepage
\end{document}
